\section*{अध्याय \ref{ch:g11:dissolution} --- आयनिक और आण्विक ठोसों का घुलना}

\begin{solution}{exo:g11:dissolution:1}
\ce{NaCl(s) -> Na+(aq) + Cl-(aq)};
\ce{CaCl2(s) -> Ca^{2+}(aq) + 2Cl-(aq)};
\ce{Na2SO4(s) -> 2Na+(aq) + SO4^{2-}(aq)};
\ce{AlCl3(s) -> Al^{3+}(aq) + 3Cl-(aq)}।
\end{solution}

\begin{solution}{exo:g11:dissolution:2}
$[\ce{Na+}] = 2 \times 0.050 = \qty{0.10}{mol/L}$;
$[\ce{SO4^{2-}}] = \qty{0.050}{mol/L}$।
\end{solution}

\begin{solution}{exo:g11:dissolution:3}
\omterm{def:g11:dissolution:solvation}{जलयोजन} \omterm{def:g9:ions:ion}{आयन} का उससे जुड़े पानी के \omterm{def:g7:atoms-and-molecules:molecule}{अणुओं} से घिर जाना है।
ऑक्सीजन वाला सिरा ($\delta^-$) \omterm{def:g9:ions:ion}{धनायन} की ओर रहता है, हाइड्रोजन वाला सिरा
($\delta^+$) \omterm{def:g9:ions:ion}{ऋणायन} की ओर: विपरीत आवेश आकर्षित करते हैं।
\end{solution}

\begin{solution}{exo:g11:dissolution:4}
$M(\ce{AlCl3}) = 27.0 + 3 \times 35.5 = \qty{133.5}{g/mol}$;
$n = 2.67 / 133.5 = \qty{0.0200}{mol}$; $c = 0.0200 / 0.2000 =
\qty{0.100}{mol/L}$; $[\ce{Al^{3+}}] = \qty{0.100}{mol/L}$,
$[\ce{Cl-}] = \qty{0.300}{mol/L}$।
\end{solution}

\begin{solution}{exo:g11:dissolution:5}
सिर \ce{-COO-} आयनिक है, पानी से आकर्षित: \omterm{def:g11:dissolution:hydrophilic}{जलरागी}। केवल \ce{C-C} और लगभग
अध्रुवीय \ce{C-H} आबंधों वाली 17 कार्बन \omterm{def:g7:atoms-and-molecules:atom}{परमाणुओं} की पूँछ
अध्रुवीय है: \omterm{def:g11:dissolution:hydrophilic}{जलविरागी}।
\end{solution}

\begin{solution}{exo:g11:dissolution:6}
नमक: पानी (आयनिक ठोस, ध्रुवीय \omterm{def:g6:solutions-and-solubility:solute}{विलायक})। मोम: साइक्लोहेक्सेन (अध्रुवीय
श्रृंखलाएँ)। चीनी: पानी (उसके \ce{O-H} समूह पानी के साथ \omterm{def:g11:polarity-and-cohesion:hydrogen-bond}{हाइड्रोजन आबंध}
बनाते हैं)। आयोडीन: साइक्लोहेक्सेन (\omterm{def:g11:polarity-and-cohesion:polar-molecule}{अध्रुवीय अणु})।
\end{solution}

\begin{solution}{exo:g11:dissolution:7}
एथेनॉल का \ce{O-H} समूह पानी के \omterm{def:g7:atoms-and-molecules:molecule}{अणुओं} के साथ \omterm{def:g11:polarity-and-cohesion:hydrogen-bond}{हाइड्रोजन आबंध} बनाता है,
जो उसके द्वारा तोड़े गए आबंधों की जगह ले लेते हैं: एथेनॉल पानी में समा जाता है।
अध्रुवीय हेक्सेन ऐसे कोई आबंध नहीं बना सकता; \omterm{def:g11:polarity-and-cohesion:hydrogen-bond}{हाइड्रोजन आबंधों} से एक-दूसरे से जुड़े
पानी के \omterm{def:g7:atoms-and-molecules:molecule}{अणु} उसे भीतर नहीं आने देते।
\end{solution}

\begin{solution}{exo:g11:dissolution:8}
\omterm{def:g11:dissolution:hydrophilic}{जलविरागी} पूँछें पानी से बचती हैं और केंद्र में चिकनाई में घुल जाती हैं;
आवेशित सिर पानी से आकर्षित होते हैं और बाहर रहते हैं। सभी
\omterm{def:g11:dissolution:amphiphilic}{मिसेलों} की सतहों पर ऋणात्मक आवेश होते हैं, और समान आवेश एक-दूसरे को
प्रतिकर्षित करते हैं: वे अलग रहते हैं।
\end{solution}

\begin{solution}{exo:g11:dissolution:9}
साइक्लोहेक्सेन: वह पानी के साथ \omterm{def:g6:solutions-and-solubility:miscible}{अमिश्रणीय} है, इसलिए एक अलग परत बनाता है जिसे
निकाला जा सकता है, और सुगंध (अध्रुवीय) उसमें अच्छी तरह \omterm{def:g2:mixing-and-dissolving:dissolve}{घुलती} है।
एथेनॉल पानी के साथ \omterm{def:g6:solutions-and-solubility:miscible}{मिश्रणीय} है: कोई दूसरी परत नहीं बनती, और कुछ भी
अलग नहीं किया जा सकता।
\end{solution}

\begin{solution}{exo:g11:dissolution:10}
साइक्लोहेक्सेन पानी से कम घना है। निचली, जलीय परत टोंटी से
निकाली जाती है। आयोडीन ऊपरी साइक्लोहेक्सेन परत में है, जो
कीप में रहती है।
\end{solution}

\begin{solution}{exo:g11:dissolution:11}
कुछ नहीं: कॉपर सल्फ़ेट आयनिक है और अध्रुवीय साइक्लोहेक्सेन में नहीं
\omterm{def:g2:mixing-and-dissolving:dissolve}{घुलता}। पानी नीला रहता है और साइक्लोहेक्सेन रंगहीन।
\end{solution}

\begin{solution}{exo:g11:dissolution:12}
$[\ce{Cl-}] = 2c$, इसलिए $c = \qty{0.150}{mol/L}$;
$n = 0.150 \times 0.2500 = \qty{0.0375}{mol}$;
$m = 0.0375 \times 111.1 = \qty{4.17}{g}$।
\end{solution}

\begin{solution}{exo:g11:dissolution:13}
कैल्सियम और मैग्नीशियम \omterm{def:g9:ions:ion}{आयन} स्टीयरेट \omterm{def:g9:ions:ion}{आयनों} को मैल के रूप में अवक्षेपित कर देते हैं:
\omterm{def:g11:dissolution:amphiphilic}{मिसेल} बनाने से पहले ही साबुन खप जाता है, और ठीक से झाग नहीं देता।
अकेला कैल्सियम: \ce{Ca^{2+}} का \qty{0.010}{mol}
स्टीयरेट का \qty{0.020}{mol} ले लेता है, यानी प्रति लीटर $0.020 \times 306.0 =
\qty{6.1}{g}$ सोडियम स्टीयरेट।
\end{solution}

\begin{solution}{exo:g11:dissolution:14}
आयतन \qty{0.2000}{L}। \ce{Na+}: $0.020 / 0.2000 = \qty{0.10}{mol/L}$;
\ce{Ca^{2+}}: $0.010 / 0.2000 = \qty{0.050}{mol/L}$; \ce{Cl-}:
$(0.020 + 0.020) / 0.2000 = \qty{0.20}{mol/L}$। धनात्मक आवेश:
$0.10 + 2 \times 0.050 = 0.20$; ऋणात्मक: $0.20$। उदासीन।
\end{solution}

\begin{solution}{exo:g11:dissolution:15}
एक लीटर हेप्टेन का द्रव्यमान \qty{0.68}{kg} है और वह
$17 \times 0.68 = \qty{12}{g}$ आयोडीन रख सकता है, जो एक लीटर पानी से लगभग 40 गुना
अधिक है। साथ हिलाने पर लगभग सारा आयोडीन हेप्टेन में चला जाता है:
\omterm{def:g10:chemical-species:extraction}{निष्कर्षण} दक्ष है।
\end{solution}

\begin{solution}{pb:g11:dissolution:1}
\textbf{1.} उन चट्टानों से जिनसे होकर पानी बहा है (चूना पत्थर,
डोलोमाइट, जिप्सम), जो थोड़ी-सी \omterm{def:g2:mixing-and-dissolving:dissolve}{घुलती} हैं।

\textbf{2.} $[\ce{Ca^{2+}}] = 0.120 / 40.1 = \qty{2.99e-3}{mol/L}$;
$[\ce{Mg^{2+}}] = 0.024 / 24.3 = \qty{9.9e-4}{mol/L}$।

\textbf{3.} \ce{Ca(HCO3)2 -> Ca^{2+}(aq) + 2HCO3-(aq)}; इसलिए
$[\ce{HCO3-}] = 2 \times \num{2.99e-3} = \qty{5.99e-3}{mol/L}$।

\textbf{4.} $\num{2.99e-3} \times 150 = \qty{0.449}{mol}$।

\textbf{5.} $18 \times 12.0 + 35 \times 1.0 + 2 \times 16.0 + 23.0 =
\qty{306.0}{g/mol}$।

\textbf{6.} \ce{C17H35COONa(s) -> C17H35COO-(aq) + Na+(aq)}।

\textbf{7.} सिर \ce{-COO-}: \omterm{def:g11:dissolution:hydrophilic}{जलरागी}; पूँछ \ce{C17H35-}:
\omterm{def:g11:dissolution:hydrophilic}{जलविरागी}। उसमें दोनों प्रकार के भाग हैं: \omterm{def:g11:dissolution:amphiphilic}{उभयरागी}।

\textbf{8.} साबुन के \omterm{def:g9:ions:ion}{आयनों} का एक नन्हा गोला, पूँछें भीतर, सिर बाहर।
पूँछें चिकनाई में घुल जाती हैं, जो \omterm{def:g11:dissolution:amphiphilic}{मिसेलों} में लिपटी बूँदों में टूट जाती
है; आवेशित सतहें बूँदों को पानी में अलग रखती हैं,
जो उन्हें बहा ले जाता है।

\textbf{9.} साबुन को \omterm{def:g9:ions:ion}{आयन} और \omterm{def:g11:dissolution:amphiphilic}{मिसेल} बनाने होते हैं, जिसके लिए सिरों के चारों ओर पानी
चाहिए; और चिकनाई धोने वाले पानी के साथ बह जाती है। तेल में चिकनाई
बस घुलकर त्वचा पर ही रह जाती।

\textbf{10.} \ce{2C17H35COO- + Ca^{2+} -> Ca(C17H35COO)2}: आवेश बाईं ओर
$2 \times (-1) + 2 = 0$, दाईं ओर $0$।

\textbf{11.} \ce{C17H35COO-}: आधिक्य $- 2x$; \ce{Ca^{2+}}:
$\num{2.99e-3} - x$; कैल्सियम स्टीयरेट: $x$। कैल्सियम \omterm{def:g9:ions:ion}{आयन}
सीमांत हैं: $x_{\max} = \qty{2.99e-3}{mol}$।

\textbf{12.} प्रति लीटर $2 \times \num{2.99e-3} = \qty{5.99e-3}{mol}$।

\textbf{13.} $M = 36 \times 12.0 + 70 \times 1.0 + 4 \times 16.0 + 40.1 =
\qty{606.1}{g/mol}$; प्रति लीटर $\num{2.99e-3} \times 606.1 = \qty{1.81}{g}$
मैल।

\textbf{14.} प्रति लीटर $2 \times \num{9.9e-4} = \qty{1.98e-3}{mol}$।

\textbf{15.} $2 \times 0.449 = \qty{0.898}{mol}$।

\textbf{16.} $0.898 \times 306.0 = \qty{275}{g}$।

\textbf{17.} मैग्नीशियम: $\num{9.9e-4} \times 150 = \qty{0.148}{mol}$,
जो स्टीयरेट का \qty{0.296}{mol}, यानी \qty{91}{g} साबुन, लेता है; कुल मिलाकर
लगभग $275 + 91 = \qty{366}{g}$।

\textbf{18.} $275 / 100 \approx 2.7$ टिकियाँ।

\textbf{19.} एक स्नान का कैल्सियम लगभग \qty{0.27}{kg} साबुन मैल के रूप में
बर्बाद करता है।
\end{solution}
